\begin{statement}{P.6}{Known}{Exact domains of the opened comparison and torsion-action theorems}
The following are the comparison statements used in this return.
They concern the full profinite groups except where a different
completion is explicitly displayed.
\begin{enumerate}[label=(\roman*)]
\item Drinfeld's equations (4.3), (4.4), (4.10) define the
two cycles and the pentagon; his construction embeds the
arithmetic action into the group of invertible solutions.
The two-cycle source is the modular free product, whereas
the pentagon is the additional five-point compatibility.
\item For $n\geq4$, Minamide--Nakamura identify
$\Out(\widehat B_n/Z(\widehat B_n))$ with
$\widehat{\mathrm{GT}}$, and consequently identify
$\Out(\widehat\Gamma_{1,2})$ with that group.
Hoshi--Minamide--Mochizuki identify the complete outer
automorphism group of the $n$-point configuration group
of the tripod, $n\geq2$, with
$\widehat{\mathrm{GT}}\times\mathfrak S_{n+3}$.
\item Boggi's 2026 v6 classification is
$\Out^{I_0}(\check\Gamma(S))\simeq\widehat{\mathrm{GT}}$
for $d(S)>1$, with the stated symmetric-group factor for
the pure group outside $S_{1,2}$. His full profinite
statement is a faithful representation into
$\Out^{I_0}(\widehat{P\Gamma}(S))$.
\item Lochak--Nakamura--Schneps construct full profinite
actions for zero or one marked point. Their Section 8
does not prove that every cutting homomorphism intertwines
these actions: preservation of the image and injectivity
of the cutting map are explicit additional requirements.
\item Lochak--Schneps identify the separate cycle equations
with nonabelian cocycle conditions for groups of orders
two, three and five. The corresponding cohomology sets
have respectively one, two and two classes. Collas proves
prime-order torsion conjugacy and cyclotomic power-conjugacy
in the full genus-zero mapping class groups; his genus-one
paper gives the analogous action for the subgroup
$I\Gamma\simeq\widehat{GR}$.
\end{enumerate}
These assertions identify neither the abelianization of
$\widehat{\mathrm{GT}}$ nor the conjugacy classes of all
involutions of $\widehat{\mathrm{GT}}_0$.
\end{statement}
\begin{proof}[Derivations and scope]
In \cite[\S4, equations (4.3), (4.4), (4.10), (4.14)]{ROOT_Drinfeld1990},
the first two relations describe the normalized modular
action and the third imposes the next braid compatibility.
The passage from the modular generators to $x,y$ is the
explicit extension--restriction calculation of P.4.
Thus dropping the last relation gives the group in this
question; the arithmetic construction in the same section
satisfies all three. No strictness assertion is needed
from that source: strictness is proved in P.2--P.4 and,
independently, B.27.

For (ii), the bar over the braid group in
\cite[Theorem B = Theorem 4.3 and Corollary C]{ROOT_MN2019}
means quotient by its center. Forgetting this bar would
lose the extra central-transvection factor in its Theorem A.
For the configuration statement, set $(g,r)=(0,3)$ in
\cite[Corollaries B and C, pp.3--4]{ROOT_HMM2017}.
Its exceptional case $(r,n)=(0,2)$ is then absent and
$n^*=n+3$. The subgroup stabilizing all generalized
fibres is exactly the pro-all-primes GT group in Corollary C.
The symmetric factor records label permutations. These
are outer-automorphism identifications of the groups
acted upon; abelianizing those source groups does not
compute the abelianization of their automorphism group.

For (iii), the check accent in
\cite[Theorem A = Theorem 6.1]{ROOT_Boggi2026}
is the congruence completion. The same source's
Theorem B and Theorem 6.7 separately supply the faithful
full profinite action. Keeping these two conclusions
separate prevents substituting a congruence isomorphism
for an unproved full-completion classification.
Neither statement evaluates a specified frame character.

For (iv), the generators and relations in
\cite[Theorem 1.1]{ROOT_LNS2026} define automorphisms
of the stated full mapping class groups. Section 8,
printed pp.19--20, discusses a cut surface with more
boundary components. Restriction to that subgroup
requires image preservation, and recovery of its action
requires injectivity of its profinite map. The authors
explicitly retain these hypotheses and explain the
known compatibility for $I\Gamma$.
Thus the existence of their action alone cannot be used
as the missing sewing identity in C.2 or as a value of
$\Delta$ on the pentagon group.

For (v), apply \cite[Proposition 3 and Theorem 2]{PBC_LS1997}:
$H^1(C_2,\widehat F_2)$ has one element,
$H^1(C_3,\widehat F_2)$ has the trivial and $xy$ classes,
and the five-cycle cohomology on $\widehat K(0,5)$ has
the two displayed classes. Theorem 2 inserts these
representatives according to $\lambda$ modulo three
and five. This computes each equation's cocycle form;
it does not assert that independently chosen finite
solutions lift to their simultaneous full intersection.
The latter passage requires compatible fibres, as
in P.3 and C.13.

Collas's genus-zero Theorems A and B
\cite{ROOT_CollasZero2012} concern torsion in
$\widehat\Gamma_{0,[n]}$. His genus-one Theorems A--D
\cite{ROOT_CollasOne2012} concern torsion in
$\widehat\Gamma_{1,[n]}$ and the action of
$I\Gamma\simeq\widehat{GR}$, whose definition contains
additional torsion-compatibility equations. In particular,
the operator group and the group containing the torsion
element are different groups. These results cannot be
read as a torsion classification of GT itself.

The same distinction applies to the 2024 sources.
Boggi--Zalesskii's Theorem A and Corollary B
\cite{ROOT_BZ2024} compare finite $p$-subgroups in a
good discrete group and its completion, with finite
virtual cohomological-type hypotheses. Boggi's
\cite[Theorem A, Corollary B and Theorem C]{ROOT_BoggiCSP2024}
treat symmetric mapping class groups and their torsion
conjugacy. No identification of the profinite operator
group in this problem as a completion satisfying those
hypotheses is supplied. Therefore these theorems do
not by themselves yield R25's torsion bound.
\end{proof}

\begin{statement}{P.7}{Known}{Finite characters, frame data, and the remaining literature distinction}
For $G=\widehat{\mathrm{GT}}$ or
$\widehat{\mathrm{GT}}_0$, let $G^{\rm ab}$ mean its
maximal continuous abelian quotient. There is a
canonical decomposition induced by the specified
arithmetic embedding and the multiplier:
\[
 G^{\rm ab}\simeq\widehat{\mathbb Z}^{\times}\times C_G,
 \qquad C_G=\ker(\lambda:G^{\rm ab}\to
                              \widehat{\mathbb Z}^{\times}).
                                                        \tag{P.7.1}
\]
Every continuous finite-order character is uniquely the
product of a character of the multiplier and a character
trivial on the specified $G_{\mathbb Q}$.
The supplied R24 gives a nonzero quadratic character of
$C_{\widehat{\mathrm{GT}}_0}$. Its restriction to
$\widehat{\mathrm{GT}}$ is not evaluated here.

The supplied, refereed R25 implies that every nontrivial
finite subgroup of either group has order two, and each
of its nontrivial elements has multiplier $-1$ at every
prime. This leaves conjugacy to $(-1,1)$ undecided.
The elliptic frame classes have antecedents as torsion
conjugators in the opened literature. The particular
quadratic sum $\Delta$, its extension to the full gentle
group, and its nonzero full lift are additional data
in R24, not conclusions of the opened torsion-action
theorems.
\end{statement}
\begin{proof}
The multiplier restricted to $G_{\mathbb Q}$ is the
cyclotomic character. Kronecker--Weber identifies the
arithmetic abelianization with
$\widehat{\mathbb Z}^{\times}$. The composite
$G_{\mathbb Q}\to G\to G^{\rm ab}$ therefore factors
through a homomorphism
$j:\widehat{\mathbb Z}^{\times}\to G^{\rm ab}$ with
$\lambda j=1$. Multiplication maps
$\widehat{\mathbb Z}^{\times}\times C_G$ bijectively
to $G^{\rm ab}$; its inverse sends $g$ to
$(\lambda(g),g\,j(\lambda(g))^{-1})$. Both maps are
continuous. This proves (P.7.1).
For a character $\eta$ with finite abelian target,
its restriction to arithmetic determines the first
factor; dividing by that character of $\lambda$
leaves precisely a character of $C_G$. Thus the
Dirichlet characters give all characters in the first
factor, while the second factor is exactly the extra
question. R24's $\Delta$ is trivial on arithmetic and
takes the value one on E.58, proving the stated
nontriviality. This is not a computation of $C_G$.
Statement C.23 below adds the nonzero arithmetic-trivial character
$\Delta+\Delta\circ D$ from an actual full element with two
retained evaluations. This new construction does not yet prove
independence from $\Delta$: the two characters coincide if
$\Delta\circ D$ is zero. The existence and comparison in C.23
are part of the present proof, rather than a claim about the older
literature.

The normalized modular restriction embeds both groups
in the pure peripheral outer group $\mathcal P$.
Apply R25 D.32 to a finite subgroup and D.34 to its
nontrivial element. The first makes the subgroup have
order at most two; the second determines every
component of its multiplier. This is the complete
reduction of the cited R25 input to the present groups.
P.1's self-centralizer computation applies to the
particular element $(-1,1)$ and supplies no conjugator
for another involution. The older question remains
explicit in \cite[\S1.4, questions (2), (4)]{PBC_SchnepsSurvey};
our search has not supplied a later general conjugacy
theorem.

To compare frame data precisely, a power-conjugacy
statement gives an equality
$\Phi(r)=A r^{\lambda}A^{-1}$. In the modular
free product the finite-factor centralizer theorem
reduces the ambiguity in $A$; choosing the pure
normalization gives the frame in R24. Collas's
statements assert existence of such conjugacy, and
Lochak--Schneps' cocycle representatives describe its
two- and three-cycle origins. They do not compare
the reductions of both normalized frames at twelve
or form their quadratic sum.

Schneps's special-locus theorem makes a stronger,
but explicitly conditional, comparison:
\cite[Theorem 3.1]{ROOT_SchnepsSpecial2006} identifies
two extra equations with commutation of two specific
covering/filling diagrams; its Theorem 3.4 derives
the listed harmonic and genus-one relations under
those equations. The reduction is to naturality
of the specified maps, not to torsion conjugacy alone.
Similarly \cite[Theorem 1 and Proposition 1]{ROOT_Ihara2000}
prove the Kummer filling relation for Galois and
identify it with equivariance of the map
$x^n\mapsto x$, $y\mapsto y$, other conjugate
punctures $\mapsto1$. Its further Propositions 2--3
compare the resulting subgroup with distribution
and logarithmic gamma identities; they do not
identify it with all of GT.

Finally the character constructed by Burau rigidity
in \cite[\S4]{ROOT_Marin2008} is a homomorphism on
the multiplier-one subgroup and extends as a twisted
one-cocycle on the whole group. Its values are
$\mathbb Z_\ell$-adic. Such a cocycle is not,
merely by its name, a finite-order character of
the whole profinite group. Formula (P.7.1) states
the exact remaining character question.

Within the sources actually opened, no theorem was
found giving R24's nonzero $\Delta$ or R25's full
torsion bound. This is a bounded comparison, not
a claim of historical priority. P.4 and B.27 give
actual non-Galois elements of the full gentle
group, since the Galois image satisfies the pentagon.
The original 2003 harmonic paper remains an access gap.
The full 2010 three-author paper was subsequently obtained;
its exact arithmetic comparison is recorded in P.8.
\end{proof}

\begin{statement}{P.8}{Known}{The exact arithmetic scope of the equianharmonic matrix formula}
Put
\[
 A=\begin{pmatrix}1&2\\0&1\end{pmatrix},\quad
 B=\begin{pmatrix}1&0\\-2&1\end{pmatrix},\quad
 R=\begin{pmatrix}0&1\\-1&1\end{pmatrix}.
\]
For $\sigma\in G_{\mathbb Q^{\rm ab}}$, the normalized
Lochak--Schneps equianharmonic parameter satisfies
\[
 h_\sigma(A,B)=(-1)^{\rho_3(\sigma)/2-2\rho_2(\sigma)}
 B_\sigma(R,R^2)
 \begin{pmatrix}1&-8\rho_2(\sigma)\\0&1\end{pmatrix}.
                                                        \tag{P.8.1}
\]
Here $B_\sigma$ is the Anderson--Ihara beta function.
The exponent is integral because $\sigma$ fixes $\sqrt3$.
This is an arithmetic frame formula; it is not asserted
for every profinite GT element.
\end{statement}
\begin{proof}
Set $\lambda=1$ in \cite[Theorem 4.16, p.446]{ROOT_NTY2010}
and substitute its equation (4.8). This gives (P.8.1),
retaining its central sign and Kummer shear.
The proof uses the Ferrari kernel and the Mordell
elliptic Tate module: Proposition 4.4 computes the beta
action, while Proposition 4.11 supplies the actual path
conjugator. Proposition 4.13 then gives power-conjugacy
of three specified braids for Galois elements.
Section 5 corrects the sign in the earlier braid
formula and derives its first case from equation
(II$'$). The introduction, p.432, explicitly leaves
its extension to all GT elements undecided.
Thus elliptic frame specializations and their signs
have this precise antecedent. The nonzero full gentle
parity character and the lift in R24 are additional
claims, not consequences of that arithmetic formula.
\end{proof}
